% Options for packages loaded elsewhere
% Options for packages loaded elsewhere
\PassOptionsToPackage{unicode}{hyperref}
\PassOptionsToPackage{hyphens}{url}
%
\documentclass[
  10pt,
  ignorenonframetext,
  aspectratio=169,
  spanish,
]{beamer}
\newif\ifbibliography
\usepackage{pgfpages}
\setbeamertemplate{caption}[numbered]
\setbeamertemplate{caption label separator}{: }
\setbeamercolor{caption name}{fg=normal text.fg}
\beamertemplatenavigationsymbolsempty
% remove section numbering
\setbeamertemplate{part page}{
  \centering
  \begin{beamercolorbox}[sep=16pt,center]{part title}
    \usebeamerfont{part title}\insertpart\par
  \end{beamercolorbox}
}
\setbeamertemplate{section page}{
  \centering
  \begin{beamercolorbox}[sep=12pt,center]{section title}
    \usebeamerfont{section title}\insertsection\par
  \end{beamercolorbox}
}
\setbeamertemplate{subsection page}{
  \centering
  \begin{beamercolorbox}[sep=8pt,center]{subsection title}
    \usebeamerfont{subsection title}\insertsubsection\par
  \end{beamercolorbox}
}
% Prevent slide breaks in the middle of a paragraph
\widowpenalties 1 10000
\raggedbottom
\AtBeginPart{
  \frame{\partpage}
}
\AtBeginSection{
  \ifbibliography
  \else
    \frame{\sectionpage}
  \fi
}
\AtBeginSubsection{
  \frame{\subsectionpage}
}
\usepackage{iftex}
\ifPDFTeX
  \usepackage[T1]{fontenc}
  \usepackage[utf8]{inputenc}
  \usepackage{textcomp} % provide euro and other symbols
\else % if luatex or xetex
  \usepackage{unicode-math} % this also loads fontspec
  \defaultfontfeatures{Scale=MatchLowercase}
  \defaultfontfeatures[\rmfamily]{Ligatures=TeX,Scale=1}
\fi
\usepackage{lmodern}

\usetheme[]{Madrid}
\usecolortheme[]{dolphin}
\usefonttheme[]{professionalfonts}
\usefonttheme{serif} % use mainfont rather than sansfont for slide text
\ifPDFTeX\else
  % xetex/luatex font selection
  \setmainfont[]{DejaVu Sans}
  \setsansfont[]{DejaVu Sans}
  \setmonofont[]{DejaVu Sans Mono}
\fi
% Use upquote if available, for straight quotes in verbatim environments
\IfFileExists{upquote.sty}{\usepackage{upquote}}{}
\IfFileExists{microtype.sty}{% use microtype if available
  \usepackage[]{microtype}
  \UseMicrotypeSet[protrusion]{basicmath} % disable protrusion for tt fonts
}{}
\makeatletter
\@ifundefined{KOMAClassName}{% if non-KOMA class
  \IfFileExists{parskip.sty}{%
    \usepackage{parskip}
  }{% else
    \setlength{\parindent}{0pt}
    \setlength{\parskip}{6pt plus 2pt minus 1pt}}
}{% if KOMA class
  \KOMAoptions{parskip=half}}
\makeatother


\usepackage{longtable,booktabs,array}
\usepackage{calc} % for calculating minipage widths
\usepackage{caption}
% Make caption package work with longtable
\makeatletter
\def\fnum@table{\tablename~\thetable}
\makeatother
\usepackage{graphicx}
\makeatletter
\newsavebox\pandoc@box
\newcommand*\pandocbounded[1]{% scales image to fit in text height/width
  \sbox\pandoc@box{#1}%
  \Gscale@div\@tempa{\textheight}{\dimexpr\ht\pandoc@box+\dp\pandoc@box\relax}%
  \Gscale@div\@tempb{\linewidth}{\wd\pandoc@box}%
  \ifdim\@tempb\p@<\@tempa\p@\let\@tempa\@tempb\fi% select the smaller of both
  \ifdim\@tempa\p@<\p@\scalebox{\@tempa}{\usebox\pandoc@box}%
  \else\usebox{\pandoc@box}%
  \fi%
}
% Set default figure placement to htbp
\def\fps@figure{htbp}
\makeatother


% definitions for citeproc citations
\NewDocumentCommand\citeproctext{}{}
\NewDocumentCommand\citeproc{mm}{%
  \begingroup\def\citeproctext{#2}\cite{#1}\endgroup}
\makeatletter
 % allow citations to break across lines
 \let\@cite@ofmt\@firstofone
 % avoid brackets around text for \cite:
 \def\@biblabel#1{}
 \def\@cite#1#2{{#1\if@tempswa , #2\fi}}
\makeatother
\newlength{\cslhangindent}
\setlength{\cslhangindent}{1.5em}
\newlength{\csllabelwidth}
\setlength{\csllabelwidth}{3em}
\newenvironment{CSLReferences}[2] % #1 hanging-indent, #2 entry-spacing
 {\begin{list}{}{%
  \setlength{\itemindent}{0pt}
  \setlength{\leftmargin}{0pt}
  \setlength{\parsep}{0pt}
  % turn on hanging indent if param 1 is 1
  \ifodd #1
   \setlength{\leftmargin}{\cslhangindent}
   \setlength{\itemindent}{-1\cslhangindent}
  \fi
  % set entry spacing
  \setlength{\itemsep}{#2\baselineskip}}}
 {\end{list}}
\usepackage{calc}
\newcommand{\CSLBlock}[1]{\hfill\break\parbox[t]{\linewidth}{\strut\ignorespaces#1\strut}}
\newcommand{\CSLLeftMargin}[1]{\parbox[t]{\csllabelwidth}{\strut#1\strut}}
\newcommand{\CSLRightInline}[1]{\parbox[t]{\linewidth - \csllabelwidth}{\strut#1\strut}}
\newcommand{\CSLIndent}[1]{\hspace{\cslhangindent}#1}

\ifLuaTeX
\usepackage[bidi=basic]{babel}
\else
\usepackage[bidi=default]{babel}
\fi
\ifPDFTeX
\else
\babelfont{rm}[]{DejaVu Sans}
\fi
% get rid of language-specific shorthands (see #6817):
\let\LanguageShortHands\languageshorthands
\def\languageshorthands#1{}


\setlength{\emergencystretch}{3em} % prevent overfull lines

\providecommand{\tightlist}{%
  \setlength{\itemsep}{0pt}\setlength{\parskip}{0pt}}



 


% Estilo separado del contenido. Compilación comprobada con XeLaTeX.
\definecolor{SYSBAzul}{HTML}{0B3954}
\definecolor{SYSBVerde}{HTML}{168C91}
\definecolor{SYSBTexto}{HTML}{203E50}
\setbeamercolor{structure}{fg=SYSBAzul}
\setbeamercolor{normal text}{fg=SYSBTexto,bg=white}
\setbeamercolor{frametitle}{fg=white,bg=SYSBAzul}
\setbeamercolor{block title}{fg=white,bg=SYSBAzul}
\setbeamercolor{block body}{fg=SYSBTexto,bg=SYSBAzul!4}
\setbeamerfont{frametitle}{size=\large,series=\bfseries}
\setbeamerfont{footline}{size=\tiny}
\setbeamersize{text margin left=6mm,text margin right=6mm}
\setbeamertemplate{navigation symbols}{}
\setbeamertemplate{footline}{\hfill\textcolor{SYSBAzul}{SYSB\quad\insertframenumber/\inserttotalframenumber}\hspace{6mm}\vspace{2mm}}
\setbeameroption{hide notes}
\setlength{\parskip}{2pt}
\setlength{\abovedisplayskip}{5pt}
\setlength{\belowdisplayskip}{5pt}
\hypersetup{pdftitle={Sistemas y Señales Biomédicos},pdfauthor={Pablo Eduardo Caicedo-Rodríguez}}

\makeatletter
\@ifpackageloaded{caption}{}{\usepackage{caption}}
\AtBeginDocument{%
\ifdefined\contentsname
  \renewcommand*\contentsname{Tabla de contenidos}
\else
  \newcommand\contentsname{Tabla de contenidos}
\fi
\ifdefined\listfigurename
  \renewcommand*\listfigurename{Listado de Figuras}
\else
  \newcommand\listfigurename{Listado de Figuras}
\fi
\ifdefined\listtablename
  \renewcommand*\listtablename{Listado de Tablas}
\else
  \newcommand\listtablename{Listado de Tablas}
\fi
\ifdefined\figurename
  \renewcommand*\figurename{Figura}
\else
  \newcommand\figurename{Figura}
\fi
\ifdefined\tablename
  \renewcommand*\tablename{Tabla}
\else
  \newcommand\tablename{Tabla}
\fi
}
\@ifpackageloaded{float}{}{\usepackage{float}}
\floatstyle{ruled}
\@ifundefined{c@chapter}{\newfloat{codelisting}{h}{lop}}{\newfloat{codelisting}{h}{lop}[chapter]}
\floatname{codelisting}{Listado}
\newcommand*\listoflistings{\listof{codelisting}{Listado de Listados}}
\makeatother
\makeatletter
\makeatother
\makeatletter
\@ifpackageloaded{caption}{}{\usepackage{caption}}
\@ifpackageloaded{subcaption}{}{\usepackage{subcaption}}
\makeatother

\usepackage{bookmark}
\IfFileExists{xurl.sty}{\usepackage{xurl}}{} % add URL line breaks if available
\urlstyle{same}
\hypersetup{
  pdftitle={Revisión académica de las diapositivas SYSB},
  pdflang={es},
  hidelinks,
  pdfcreator={LaTeX via pandoc}}


\title{Revisión académica de las diapositivas SYSB}
\subtitle{Microcurrículo, libro guía y correspondencia por diapositiva}
\author{}
\date{2026-09-09}

\begin{document}
\frame{\titlepage}


\begin{frame}{Dictamen}
\phantomsection\label{dictamen}
\textbf{Los tres archivos originales seguían la secuencia temática de
las 30 sesiones del microcurrículo, pero su correspondencia con el libro
guía era parcial y contenían errores de formulación y de
implementación.} El problema principal no era el orden de las semanas,
sino la ausencia de referencias verificadas, la cobertura insuficiente
de algunos fundamentos y la falta de distinción entre contenido del
libro y ampliaciones curriculares.

La revisión mantiene las 15 semanas, sus dos sesiones y el uso de
Python. Corrige los errores identificados, incorpora desarrollos
matemáticos y añade referencias en las notas de todas las diapositivas,
incluidas las portadas, separadores y síntesis. Las notas distinguen
referencias directas, adaptaciones y antecedentes de temas que el libro
guía no desarrolla.

La alineación resultante se refiere al \textbf{cronograma suministrado y
a los fundamentos pertinentes del libro}. No significa cubrir la
totalidad del libro, que incluye circuitos y otros capítulos fuera del
alcance programado, ni resolver las inconsistencias administrativas del
microcurrículo.
\end{frame}

\begin{frame}[fragile]{Fuentes efectivamente utilizadas}
\phantomsection\label{fuentes-efectivamente-utilizadas}
\begin{longtable}[]{@{}
  >{\raggedright\arraybackslash}p{(\linewidth - 2\tabcolsep) * \real{0.5000}}
  >{\raggedright\arraybackslash}p{(\linewidth - 2\tabcolsep) * \real{0.5000}}@{}}
\toprule\noalign{}
\begin{minipage}[b]{\linewidth}\raggedright
Fuente
\end{minipage} & \begin{minipage}[b]{\linewidth}\raggedright
Función en la revisión
\end{minipage} \\
\midrule\noalign{}
\endhead
\texttt{Programa\_SYSB\_2027-1.pdf} & Temas, 30 sesiones, actividades,
resultados de aprendizaje y evaluación \\
\texttt{Circuits,\ Signals\ and\ Systems\ for\ Bioengineers.pdf} & Texto
guía: John Semmlow, \textbf{3.ª edición, 2018} \\
\texttt{BIOSIGNAL\ and\ MEDICAL\ IMAGE\ PROCESSING.pdf} & Complemento de
Semmlow y Griffel, \textbf{3.ª edición, 2014}: wavelets, señal analítica
y clasificación \\
\texttt{Digital\_Signal\_Processing\_Illustration\_Using\_Python\_Springer\_2024.pdf}
& Complemento de Esakkirajan, Veerakumar y Subudhi, 2024: implementación
DSP en Python \\
Documentación oficial de SciPy, PyWavelets y scikit-learn & Convenciones
computacionales, escalamiento y condiciones de uso \\
\bottomrule\noalign{}
\end{longtable}

La revisión no atribuye a los demás libros adjuntos una validación que
no se haya realizado. Las figuras incorporadas se calculan con señales
sintéticas y no reproducen las ilustraciones del libro. Los ejemplos
nuevos son elaboraciones docentes.

\begin{block}{Discrepancias dentro del microcurrículo}
\phantomsection\label{discrepancias-dentro-del-microcurruxedculo}
\begin{longtable}[]{@{}
  >{\raggedright\arraybackslash}p{(\linewidth - 4\tabcolsep) * \real{0.3333}}
  >{\raggedright\arraybackslash}p{(\linewidth - 4\tabcolsep) * \real{0.3333}}
  >{\raggedright\arraybackslash}p{(\linewidth - 4\tabcolsep) * \real{0.3333}}@{}}
\toprule\noalign{}
\begin{minipage}[b]{\linewidth}\raggedright
Hallazgo
\end{minipage} & \begin{minipage}[b]{\linewidth}\raggedright
Evidencia
\end{minipage} & \begin{minipage}[b]{\linewidth}\raggedright
Decisión adoptada
\end{minipage} \\
\midrule\noalign{}
\endhead
Nombre de archivo e identificación interna diferentes & El nombre dice
2027-1; portada y encabezados dicen 2022-2 & Registrar la
inconsistencia. No cambiar el período institucional por inferencia \\
Edición bibliográfica diferente & El programa cita \emph{Signals and
Systems for Bioengineers}, 2.ª ed., 2012; el PDF guía suministrado es la
3.ª ed., 2018, con otro título y organización & Usar las secciones
reales de la \textbf{edición aportada} \\
Numeración del cronograma diferente de la del libro & Por ejemplo, el
cronograma sitúa sinusoides en 2.2, Laplace en 5.6--5.7 y wavelets en
7.4--7.7 & Mantener esos números como identificadores curriculares y
construir una correspondencia separada \\
Horas dirigidas incongruentes & Información general: 4.5 h; cronograma:
dos sesiones de 90 min, es decir, 3 h por semana & Conservar el
cronograma solicitado. La distribución de las 1.5 h restantes requiere
definición institucional \\
Denominación del 16 \% final ambigua & Lista general: Proyecto Final 16
\%; tabla del tercer corte: Laboratorios finales 16 \% & No inventar una
equivalencia. Remitir los porcentajes al programa que confirme el
profesor \\
\bottomrule\noalign{}
\end{longtable}

Estas discrepancias pertenecen al documento de referencia y no pueden
solucionarse honestamente cambiando las diapositivas para que aparenten
concordancia.
\end{block}
\end{frame}

\begin{frame}{Correspondencia entre las sesiones y el libro guía}
\phantomsection\label{correspondencia-entre-las-sesiones-y-el-libro-guuxeda}
Los números siguientes son \textbf{secciones reales de la 3.ª edición}.
Las páginas incluidas en las notas son posiciones del visor del PDF
aportado, no páginas impresas. Se contrastaron con los encabezados del
cuerpo del libro, porque algunos destinos de sus marcadores internos son
incorrectos.

\begin{longtable}[]{@{}
  >{\raggedright\arraybackslash}p{(\linewidth - 4\tabcolsep) * \real{0.3333}}
  >{\raggedright\arraybackslash}p{(\linewidth - 4\tabcolsep) * \real{0.3333}}
  >{\raggedright\arraybackslash}p{(\linewidth - 4\tabcolsep) * \real{0.3333}}@{}}
\toprule\noalign{}
\begin{minipage}[b]{\linewidth}\raggedright
Semana / sesión
\end{minipage} & \begin{minipage}[b]{\linewidth}\raggedright
Tema del microcurrículo
\end{minipage} & \begin{minipage}[b]{\linewidth}\raggedright
Secciones del libro guía y alcance
\end{minipage} \\
\midrule\noalign{}
\endhead
1 / 1 & Señal, sistema y bioseñal & 1.1--1.2; 1.4--1.4.5 \\
1 / 2 & Origen fisiológico, ruido y artefactos & 1.2.1--1.2.2;
1.3.1--1.3.2 \\
2 / 1 & Representación y descriptores temporales & 1.2.3; 2.2.1--2.2.2.
Se añade promediado de realizaciones: 2.2.3--2.2.4 \\
2 / 2 & Sinusoides & \textbf{2.3.1--2.3.2}, no 2.2 \\
3 / 1 & Operaciones temporales & Antecedentes: 1.4.2 y 2.4.4; desarrollo
docente específico de composición de operaciones \\
3 / 2 & Envolvente, cruces, eventos & Ampliación curricular con
antecedentes en 2.2 y 2.4. Se añaden correlación y correlación cruzada:
2.4.1, 2.4.3--2.4.4 \\
4 / 1 & Segmentación y normalización & Aplicación de 2.2.4 y 2.4.3; se
añade autocorrelación/autocovarianza: 2.4.5--2.4.6 \\
4 / 2 & Relación tiempo--frecuencia & 3.2 y 3.3.1; puente de
ortogonalidad: 2.4.2. Estacionariedad: 10.2, como anticipación \\
5 / 1 & Series de Fourier & 3.3.2--3.3.6 \\
5 / 2 & Fourier continuo, magnitud y fase & 3.3.7 y 3.4.3; propiedades y
distribuciones como ampliación formal \\
6 / 1 & Propiedades de Fourier & 3.3.7 y 5.6; integración/derivación:
6.5.2--6.5.3 y 7.2.2 \\
6 / 2 & DFT y FFT & 3.4.1--3.4.4 \\
7 / 1 & Resolución, fuga y ventanas & 3.4.3; 4.2.3--4.2.4 \\
7 / 2 & Adquisición y almacenamiento & 1.2.2--1.2.3; 4.2.1--4.2.2.
Metadatos como ampliación reproducible \\
8 / 1 & Muestreo y aliasing & 4.2.1; 5.6.1 \\
8 / 2 & Potencia, PSD y periodograma & 4.3. Enlace con autocorrelación:
2.4.5--2.4.6 \\
9 / 1 & Welch y ancho de banda & \textbf{4.4--4.5}, no solo 4.3 \\
9 / 2 & No estacionariedad y STFT & 4.6 y 10.2 \\
10 / 1 & Ventanas tiempo--frecuencia & 4.2.4 y 4.6 \\
10 / 2 & Linealidad, invariancia e impulso & \textbf{5.2--5.4};
causalidad: 1.4.3; frecuencia: 6.3 y 6.7 \\
11 / 1 & Convolución & \textbf{5.4--5.4.2} \\
11 / 2 & Suavizadores y diferenciadores & \textbf{5.5} y 8.5.1 \\
12 / 1 & Convolución en frecuencia & 5.6 y 6.4--6.7. Se incorporan
modelos, fasores y Bode: 6.2--6.6 \\
12 / 2 & Laplace, ROC, polos y ceros & \textbf{7.2--7.3}. ROC bilateral
como ampliación explícita \\
13 / 1 & Condiciones iniciales e identificación & \textbf{7.4--7.6} \\
13 / 2 & Z, diferencias y FIR & \textbf{8.2--8.5}, más 8.6.1 para
implementación \\
14 / 1 & IIR y comparación de filtros & \textbf{8.4.1--8.4.3 y 8.6.2} \\
14 / 2 & Wavelets & \textbf{Sin sección directa en el libro guía.}
Complemento: Semmlow y Griffel, 7.1--7.2; Python: Esakkirajan et al.,
5.7 \\
15 / 1 & Multirresolución & \textbf{Sin sección directa.} Complemento:
Semmlow y Griffel, 7.3; Python: Esakkirajan et al., 5.8 \\
15 / 2 & Sistemas lineales para clasificación & \textbf{Sin sección
directa para clasificación.} Antecedentes: 1.4.4 y 8.4. Complemento:
Semmlow y Griffel, 16.2--16.3 y 16.6 \\
\bottomrule\noalign{}
\end{longtable}

\begin{block}{Qué faltaba respecto del libro}
\phantomsection\label{quuxe9-faltaba-respecto-del-libro}
El capítulo 2 no es solamente una colección de descriptores: utiliza
correlación, ortogonalidad, correlación cruzada y autocorrelación como
puente hacia Fourier. Su cobertura era insuficiente en las diapositivas.
Se incorporan definiciones, diferencias entre medidas, un ejemplo de
retardo y el vínculo estadístico con la PSD.

El capítulo 6 desarrolla modelos de sistemas, fasores, elementos
básicos, Bode y combinación de transferencias. La exposición original de
semanas 11--15 saltaba demasiado pronto de una respuesta frecuencial
genérica a Laplace. Se añaden estas conexiones sin cambiar las sesiones
del cronograma.

El desarrollo original de Laplace y Z enunciaba definiciones y proponía
ejercicios, pero no mostraba suficientes pasos intermedios. Se añaden
convolución por intervalos, uso correcto del factor temporal en una
suma, respuesta con estado inicial, contraejemplo del valor final y una
ecuación en diferencias resuelta hasta su respuesta al impulso.
\end{block}
\end{frame}

\begin{frame}[fragile]{Correcciones teóricas y computacionales}
\phantomsection\label{correcciones-teuxf3ricas-y-computacionales}
\begin{longtable}[]{@{}
  >{\raggedright\arraybackslash}p{(\linewidth - 4\tabcolsep) * \real{0.3333}}
  >{\raggedright\arraybackslash}p{(\linewidth - 4\tabcolsep) * \real{0.3333}}
  >{\raggedright\arraybackslash}p{(\linewidth - 4\tabcolsep) * \real{0.3333}}@{}}
\toprule\noalign{}
\begin{minipage}[b]{\linewidth}\raggedright
Tema
\end{minipage} & \begin{minipage}[b]{\linewidth}\raggedright
Problema original
\end{minipage} & \begin{minipage}[b]{\linewidth}\raggedright
Corrección
\end{minipage} \\
\midrule\noalign{}
\endhead
RMS & Exponente escrito como \texttt{\textbackslash{}bar\{x\}\^{}\{,2\}}
& \(x_{RMS}^2=\sigma_x^2+\bar{x}^2\), usando denominador \(N\) \\
Varianza & Denominación poblacional sin distinguir estimadores &
Separación entre descripción del segmento con \(N\) y estimación con
\(N-1\) \\
Energía & Posible confusión entre suma discreta e integral física &
Explicitar \(T_s\sum|x[n]|^2\) y sus unidades \\
Transformaciones temporales & Se afirmaba que desplazar por \(b\) y
después escalar no producía \(x(at-b)\) & Mostrar ambos órdenes
equivalentes y la transformación de un evento \(t_k\mapsto(t_k+b)/a\) \\
Detector de eventos & Matriz binaria con verdaderos negativos sin
definir oportunidades negativas & Emparejamiento uno a uno, tolerancia,
VP/FP/FN y error temporal \\
Envolvente de Hilbert & Podía leerse como intensidad fisiológica general
& Condiciones de banda, separación de escalas y advertencia sobre
batimientos \\
Normalización & No se explicitaban segmentos constantes ni diferencia
entre parámetros globales y locales & Denominadores no nulos y alcance
de la regla de ajustar solo con entrenamiento \\
Fourier y tiempo & Se sugería pérdida de localización como propiedad de
toda transformada & Distinguir transformada compleja invertible de
resúmenes de magnitud o PSD \\
Sinusoides y LTI & Frecuencia conservada sin indicar existencia de
respuesta ni transitorio & Condiciones de convergencia y régimen
permanente \\
Integración & Término en DC dejado indeterminado & Valor principal y
término distribucional explícitos \\
DFT & Eje manual no coincidente con \texttt{fftfreq} en Nyquist; un
fragmento no trataba \(N\) impar & Convención exacta y pruebas de
amplitud con ambas paridades \\
Amplitud & Semiespectro bilateral confundible con amplitud unilateral &
Etiquetas correctas, duplicación de bins y supuestos de tono alineado \\
Parseval con ventana & Comparación sugerida con potencia temporal sin la
misma ponderación & Comparar con \(\sum|wx_0|^2/\sum w^2\); sumar bins
completos para la prueba exacta \\
Ancho de banda & Media y mediana espectral listadas como definiciones de
ancho & Distinguir extensión de localización del espectro \\
Hamming & Extremos no nulos presentados como discontinuidad periódica
inevitable & Diferenciar extremos de la ventana y extremos de la señal
ventaneada \\
Muestreo & Un margen de 5 Hz se declaraba imposible para cualquier
filtro real & Pedir atenuación, orden y tolerancia antes de concluir
factibilidad \\
Incertidumbre & Constante \(C\) no definida & Definir dispersión de
energía y \(\sigma_t\sigma_f\ge1/(4\pi)\) \\
STFT & Módulo cuadrado sin normalización cuantitativa completa &
\texttt{scaling="psd"} y combinación explícita de potencia unilateral \\
Promedio móvil & Implementación \texttt{same} usada junto con
interpretación de retardo causal & \texttt{lfilter}, estado inicial y
comportamiento en bordes \\
Diferenciación & Diferencia presentada sin escala temporal & Explicitar
división por \(T_s\) cuando se aproxima derivada física \\
Laplace & Mezcla de bilateral y unilateral; polos descritos solo en
rad/s & Separar transformadas; dimensión de \(s\) en s\(^{-1}\) y
significado de sus componentes \\
Modelo RC & Ganancia unitaria atribuida a membrana sin definir entrada &
Separar voltaje--voltaje de corriente--voltaje, esta última con ganancia
en ohmios \\
Estabilidad & Regla de polos sin aclarar transferencia propia y término
directo & Precisar condiciones de la regla y caso distribucional de
ganancia directa \\
Valor inicial/final & Condiciones descritas de forma conjunta & Separar
hipótesis e incorporar contraejemplo oscilatorio \\
FIR & Cortes nominales confundibles con banda plana & Especificar
bordes, rizado, atenuación y convención de \texttt{firwin} \\
Adelante--atrás & Se describía la eliminación de fase sin su cambio de
magnitud & Respuesta ideal \(|H|^2\) y evaluación de la respuesta
efectivamente aplicada \\
CWT & Escalas 1--2 con Morlet compleja podían introducir aliasing &
Selección desde pseudofrecuencias, margen y revisión de escalas
pequeñas \\
DWT & Faltaban control de nivel, orden de coeficientes y reconstrucción
de bandas & Verificación de longitud, nivel máximo y síntesis a la tasa
original \\
Clasificación & Logaritmo de una cantidad dimensional y partición solo
declarada & Cociente adimensional y ejemplo explícito de división por
sujeto \\
\bottomrule\noalign{}
\end{longtable}

La concordancia con un texto no justifica repetir una afirmación débil.
Por ejemplo, la equivalencia entre ortogonalidad y correlación centrada
requiere precisar el centrado; las diapositivas hacen esa distinción.
Del mismo modo, que el libro privilegie Laplace unilateral no convierte
la bilateral en incorrecta: se conserva como ampliación claramente
identificada.
\end{frame}

\begin{frame}{Ajustes pedagógicos y didácticos}
\phantomsection\label{ajustes-pedaguxf3gicos-y-diduxe1cticos}
Se conservan las preguntas biomédicas y los criterios de interpretación
ya útiles. En las semanas 11--15 se sustituyen las preguntas repetidas
sobre ``el parámetro principal'' por comprobaciones concretas de
soporte, retardo, polos, condiciones de los teoremas, frecuencia física,
reconstrucción y unidad independiente.

Se propone un uso de 90 minutos con apertura, fundamento, aplicación en
parejas y comprobación individual. Es una pauta docente, no una medición
del tiempo que tardará cualquier grupo. La cantidad de diapositivas
responde a su doble uso como apoyo de clase y material de estudio:
\textbf{no se recomienda leerlas todas de forma expositiva}.
Correlación, Bode y desarrollos agregados deben seleccionarse según
conocimientos previos y acompañarse de la lectura indicada.

Los RAE no se satisfacen solo por incluir temas. Se añaden evidencias
observables de lectura en inglés, explicación oral, cálculo individual,
programación reproducible, escritura científica y colaboración. Las
actividades integradoras requieren procedimiento, unidades, comprobación
e interpretación. No se altera el esquema institucional de porcentajes
ni se inventan fechas de exámenes.

La ampliación de clasificación conserva la distinción entre un banco de
filtros LTI, características no lineales y una frontera de decisión
afín. El ejemplo LDA generativo se respalda adicionalmente con la
documentación de scikit-learn: el capítulo 16 del complemento introduce
discriminadores lineales, pero no se le atribuye esa derivación exacta.
\end{frame}

\begin{frame}[fragile]{Notas y trazabilidad}
\phantomsection\label{notas-y-trazabilidad}
Cada diapositiva contiene un bloque \texttt{:::\ \{.notes\}} con el
libro y edición, número y título de sección, página del visor cuando
corresponde, clase de correspondencia y sesión curricular. Los temas
ausentes del libro se señalan como tales y se remiten al complemento
adecuado.

Puede consultar las notas en la vista del presentador de Revealjs o
leerlas de forma continua en \texttt{notas-del-docente.html}. Los PDF
Beamer se entregan con las notas ocultas para proyección; los bloques
permanecen en el QMD y en el TEX generado. El archivo
\texttt{correspondencia-diapositivas.json} permite revisar o reutilizar
el mapeo completo.

Se reemplazan las rutas externas que no venían adjuntas por recursos
incluidos en el paquete. El formato continúa separado del contenido en
\texttt{\_metadata.yml}, \texttt{estilos/reveal.scss},
\texttt{estilos/beamer-header.tex} y un filtro que representa los
callouts como bloques nativos de Beamer. La apariencia se basa en los
colores y proporción disponibles en los QMD; \textbf{no puede afirmarse
identidad visual con los estilos originales que no fueron aportados}. El
logo externo tampoco estaba disponible.
\end{frame}

\begin{frame}{Comprobaciones de la entrega}
\phantomsection\label{comprobaciones-de-la-entrega}
\begin{longtable}[]{@{}lrr@{}}
\toprule\noalign{}
Archivo & Diapositivas & Bloques de notas \\
\midrule\noalign{}
\endhead
Semanas 01--05 & 109 & 109 \\
Semanas 06--10 & 172 & 172 \\
Semanas 11--15 & 162 & 162 \\
\textbf{Total} & \textbf{443} & \textbf{443} \\
\bottomrule\noalign{}
\end{longtable}

Los tres QMD compilan a Revealjs y Beamer con Quarto 1.7.34. Se comprobó
la presencia de las notas en el HTML y en el TEX, la sintaxis de los 29
fragmentos Python y la compatibilidad de las expresiones matemáticas con
el motor KaTeX utilizado por los HTML.

Los casos numéricos reproducibles comprueban la identidad RMS, la
amplitud unilateral con número par e impar de muestras, Parseval con la
misma ventana, un retardo conocido, la aproximación de una convolución
continua, el corte de un primer orden y la reconstrucción DWT. El error
de Parseval fue aproximadamente \(5.6\times10^{-16}\); el error relativo
de reconstrucción DWT fue \(3.9\times10^{-16}\). La aproximación de
convolución usa \(T_s=0.001\) s y tiene error máximo cercano a 0.001.
Estas pruebas cubren los casos incluidos en el script; no implican haber
ejecutado todas las actividades con datos biomédicos reales.

Se revisaron visualmente páginas representativas de los PDF y se
corrigieron las alertas de texto desbordado detectadas al recorrer todas
las páginas. \textbf{La revisión visual y de navegación de los HTML
quedó limitada por el bloqueo de acceso del navegador a los archivos
locales.} Su compilación, estructura, notas y sintaxis matemática sí se
comprobaron. Las ecuaciones HTML requieren acceso a la CDN de KaTeX; los
PDF no tienen esa dependencia.
\end{frame}

\begin{frame}{Documentación computacional consultada}
\phantomsection\label{documentaciuxf3n-computacional-consultada}
\begin{itemize}
\tightlist
\item
  Normalización y potencia:
  \href{https://docs.scipy.org/doc/scipy/reference/generated/scipy.signal.periodogram.html}{SciPy,
  periodogram}.
\item
  Espectrograma:
  \href{https://docs.scipy.org/doc/scipy/reference/generated/scipy.signal.stft.html}{SciPy,
  stft}. Se conserva esta API para continuidad y se identifica su
  condición de legacy.
\item
  Diseño de filtros:
  \href{https://docs.scipy.org/doc/scipy/reference/generated/scipy.signal.firwin.html}{SciPy,
  firwin} y
  \href{https://docs.scipy.org/doc/scipy/reference/generated/scipy.signal.filtfilt.html}{filtfilt}.
\item
  Escalas y aliasing:
  \href{https://pywavelets.readthedocs.io/en/latest/ref/cwt.html}{PyWavelets,
  CWT}.
\item
  Supuestos de LDA:
  \href{https://scikit-learn.org/stable/modules/lda_qda.html}{scikit-learn,
  Linear and Quadratic Discriminant Analysis}.
\item
  Vista del presentador:
  \href{https://quarto.org/docs/presentations/revealjs/presenting.html}{Quarto,
  Presenting Slides}.
\end{itemize}
\end{frame}

\begin{frame}{Bibliografía}
\phantomsection\label{bibliografuxeda}
Texto principal: J. Semmlow (\citeproc{ref-semmlow2018}{2018}).
Complementos: J. L. Semmlow y Griffel
(\citeproc{ref-semmlowgriffel2014}{2014}) y Esakkirajan, Veerakumar, y
Subudhi (\citeproc{ref-esakkirajan2024}{2024}). El DOI de este último se
verificó en la página de créditos del PDF aportado.

\phantomsection\label{refs}
\begin{CSLReferences}{1}{0}
\bibitem[\citeproctext]{ref-esakkirajan2024}
Esakkirajan, S., T. Veerakumar, y Badri N. Subudhi. 2024. \emph{Digital
Signal Processing: Illustration Using Python}. Springer.
\url{https://doi.org/10.1007/978-981-99-6752-0}.

\bibitem[\citeproctext]{ref-semmlow2018}
Semmlow, John. 2018. \emph{Circuits, Signals, and Systems for
Bioengineers: A MATLAB-Based Introduction}. 3.ª ed. Academic Press.

\bibitem[\citeproctext]{ref-semmlowgriffel2014}
Semmlow, John L., y Benjamin Griffel. 2014. \emph{Biosignal and Medical
Image Processing}. 3.ª ed. CRC Press.

\end{CSLReferences}
\end{frame}




\end{document}
